% Reconstructed from made/mentor07.py; currently commented out in fa26/mentor07.
\begin{blocksection}
\question Write \lstinline{link_to_tree}, which takes a non-empty linked list and returns a tree in which each node has the next element as its only branch. Then write \lstinline{tree_to_link}, which does the reverse.

\begin{lstlisting}
def link_to_tree[T](lnk: Link[T]) -> Tree[T]:
    """
    >>> print(link_to_tree(Link(6, Link(2, Link(9)))))
    6
      2
        9
    """
    if ______________________________________:

        return ______________________________

    return ____________________________________________________


def tree_to_link[T](t: Tree[T]) -> Link[T]:
    """Assume every node of t has at most one branch.
    >>> print(tree_to_link(Tree(6, [Tree(2, [Tree(9)])])))
    (6 2 9)
    """
    if ______________________________________:

        return ______________________________

    return ____________________________________________________
\end{lstlisting}

\begin{solution}[1in]
\begin{lstlisting}
def link_to_tree[T](lnk: Link[T]) -> Tree[T]:
    if not isinstance(lnk.rest, Link):
        return Tree(lnk.first)
    return Tree(lnk.first, [link_to_tree(lnk.rest)])

def tree_to_link[T](t: Tree[T]) -> Link[T]:
    if is_leaf(t):
        return Link(t.label)
    return Link(t.label, tree_to_link(t.branches[0]))
\end{lstlisting}
\end{solution}

\end{blocksection}
